%% Version 1.1 (16-02-2026) https://github.com/sandervansinte/ECU-LaTeX-Template.git

% This template was designed and written in accordance with the ECU Higher Degree by Research (HDR) Thesis Preparation Guidelines (https://intranet.ecu.edu.au/research/higher-degree-by-research/resources-and-development/writing-and-presenting-your-research/thesis-formatting-for-examination).

% Nevertheless, this template was also designed with the goal of being adaptable to other institutions. So please feel free to use and modify this template for your own needs. If you have any questions or suggestions for improvement, please do not hesitate to contact me (see contact information below) or submit an issue on the GitHub repository.


% Contact: Sander Jordi van Sintemaartensdijk, MSc (Edith Cowan University, Australia)
% <s.vansintemaartensdijk@ecu.edu.au>
% https://github.com/sandervansinte


% -------------------------------------------------------------------------
%% Accreditation %%

% This template is based on the work of from various sources and the collaborative effort of other and my own. Thus I want to accredit the following people for their indirect contributions to this template:

% TU Wien ACE LaTeX Thesis Template (https://www.tuwien.at/en/etit/ict/ressources)
% ing. Herbert van Sintemaartensdijk, MA (TU Wien, Austria)
% Axel Jantsch (TU Wien, Austria)

% Masters/Doctoral Thesis LaTeX Template (Version 2.5 (27/8/17)) http://www.LaTeXTemplates.com
% Steve Gunn (http://users.ecs.soton.ac.uk/srg/softwaretools/document/templates/)
% Sunil Patel (http://www.sunilpatel.co.uk/thesis-template/)


% -------------------------------------------------------------------------
%% Instructions %%

% Before startting make sure to edit the metadata.tex file (1.Metadata/metadata.tex) and setup.tex file (1.Metadata/setup.tex) according to your specific requirements.

% Customize the packages you want to use in the packages.tex file (1.Metadata/packages.tex).

% Make sure to compile all your acronyms and definitions in the acronyms-definitions.tex file (1.Metadata/acronyms-definitions.tex) if you intend to use the glossary functionality. If you do not intend to use the glossary functionality, you can simply leave the acronyms-definitions.tex file empty.  

% The template is structured in a way that the main content of your document should be included in the Mainmatter folder (3.Mainmatter/). You can create as many files as you need for the different sections and include them in the main.tex file while making sure that the custom document is selected. Additionally you'll need to comment out includepreset and uncomment the desired files in the main.

% The Figures folder (6.Figures/) is meant to store all your figures. You can set the path for the figures in the setup.tex file (1.Metadata/setup.tex).


% -------------------------------------------------------------------------
%% Update History %%

% Update Version 1.1 (26-02-2026):
% - Added the function to have preset layouts for the content that is loaded depending on the document type (article, custom, report, research proposal, thesis).
% The custom document type is the default type and it requires to just requires uncommenting the lines of content you want to include and and commenting out the includepreset function in the main.tex file.
% - Selecting a document type will now result in custom formatting and titlepages. For example, articles will uniquely not have a title page, also, the sections will not be numbered, and the document will be formatted in a double column style. 